% Generated from tex/chapters/ch04/08_構築技術.tex for the SWEBOK structure tree
\subsubsection{文法に基づく入力処理}

文法に基づく入力処理は、入力を字句や構文の規則に従って解析する技法である。代表的な例は、プログラミング言語、設定ファイル、数式、検索式、通信プロトコルの解析である。

構文解析では、入力のトークン列から解析木または構文木と呼ばれるデータ構造を作る。解析木は、その入力がどのような構造を持つかを表す。構文解析の後、意味検査やシンボル表の確認を行い、計算処理や変換処理につなげる。

文法に基づく入力処理は、入力検証にも重要である。曖昧な文字列処理で入力を解釈すると、誤解析やセキュリティ問題につながる可能性がある。
